% A fixture for tests/test_main_nerd_pages.py: a stand-in for the plain-language edition of a paper, in the macro
% vocabulary of docs/papers/01_curled_torus/paper_plain_language.tex. It is not a paper and nothing in it is a result.
% The blog's build (scripts/build_main_nerd_pages.py) is pointed at it as a post with the plain-language edition
% only, so that the path every such post takes is tested without depending on a paper that is still being written.
\documentclass[11pt,letterpaper]{article}
\usepackage[margin=0.87in]{geometry}
\usepackage[T1]{fontenc}
\usepackage{lmodern}
\usepackage{microtype}
\usepackage{amsmath}
\usepackage{xcolor}
\usepackage{graphicx}
\usepackage{tikz}
\usepackage{pgfplots}
\pgfplotsset{compat=1.18}
\usetikzlibrary{arrows.meta,calc,positioning,patterns,decorations.pathmorphing,shapes.misc}
\usepackage{booktabs,tabularx,array}
\usepackage{float}
\usepackage{needspace}
\usepackage[font=small,labelfont=bf,labelsep=period,skip=4pt]{caption}
\usepackage{fancyhdr}
\usepackage[colorlinks=true,linkcolor=teal!65!black,urlcolor=teal!65!black,citecolor=teal!65!black]{hyperref}

% Palette
\definecolor{ink}{HTML}{17324D}
\definecolor{accent}{HTML}{176B70}
\definecolor{pale}{HTML}{EDF5F4}
\definecolor{curl}{HTML}{C8742B}
\definecolor{curlpale}{HTML}{F6E3D3}
\definecolor{flat}{HTML}{2E7D8C}
\definecolor{flatpale}{HTML}{D6ECF0}
\definecolor{melt}{HTML}{8E6C8A}
\definecolor{meltpale}{HTML}{E9DEE8}
\definecolor{sun}{HTML}{E0A526}
\definecolor{no}{HTML}{B03A2E}
\definecolor{kExact}{HTML}{176B70}
\definecolor{kMeas}{HTML}{2F5F9E}
\definecolor{kPub}{HTML}{5A5A5A}
\definecolor{kOurs}{HTML}{9A5B00}
\definecolor{kClaim}{HTML}{8A2D5A}

\setlength{\parindent}{0pt}
\setlength{\parskip}{0.62em}
\renewcommand{\baselinestretch}{1.08}
\setlength{\headheight}{14pt}
\pagestyle{fancy}
\fancyhf{}
\fancyhead[L]{\small\color{ink}A stub post}
\fancyhead[R]{\small Plain-language edition}
\fancyfoot[C]{\small\thepage}
\renewcommand{\headrulewidth}{0.3pt}
\hypersetup{pdftitle={A stub post: a fixture for the tests},pdfauthor={A fixture}}

% Chapter opener, sub-heading, take-away box
\newcommand{\chapterpage}[2]{\par\Needspace{16\baselineskip}\vspace{1.4em}{\color{accent}\small\textbf{#1}}\par\vspace{0.2em}{\color{ink}\LARGE\bfseries #2}\par\vspace{0.5em}}
\newcommand{\takeaway}[1]{\par\medskip\noindent\colorbox{pale}{\parbox{\dimexpr\linewidth-2\fboxsep\relax}{\vspace{0.35em}\textbf{The point to keep:} #1\vspace{0.35em}}}\par\medskip}
\newcommand{\whyhere}[1]{\par\medskip\noindent{\color{accent}\small\textbf{Why we went here.}} {\small #1}\par\medskip}
\newcommand{\sub}[1]{\par\medskip{\color{ink}\large\bfseries #1}\par}

% Kinds of statement, as small tags
\newcommand{\tagbox}[2]{{\sffamily\scriptsize\colorbox{#1}{\textcolor{white}{\strut #2}}}}
\newcommand{\Exact}{\tagbox{kExact}{EXACT}}
\newcommand{\Meas}{\tagbox{kMeas}{MEASURED}}
\newcommand{\MeasX}{\tagbox{kMeas}{MEASURED, EXPLORATORY}}
\newcommand{\Pub}{\tagbox{kPub}{PUBLISHED}}
\newcommand{\Ours}{\tagbox{kOurs}{OUR READING, UNVERIFIED}}
\newcommand{\Claim}{\tagbox{kClaim}{THE AUTHOR'S CLAIM}}
\newcommand{\Sketch}{\tagbox{kPub}{SKETCH, NOT DATA}}

\begin{document}
\thispagestyle{empty}
{\color{accent}\small\textbf{A FIXTURE FOR THE TESTS}}\par
\vspace{1em}
{\color{ink}\fontsize{29}{33}\selectfont\bfseries A stub post}\par
\vspace{0.5em}
{\raggedright\Large Two short chapters,\\and nothing to learn from them\par}
\vspace{0.7em}
\textbf{A fixture}\quad Not a paper\par
Stands in for the plain-language edition of a technical paper~\cite{tech}
\vspace{0.9em}

This file is a fixture for the tests of the blog's build. It is not a paper, and nothing in it is a
result. It holds one or two of each construct the build converts.

It cites one source~\cite{first} and the technical paper it stands beside~\cite{tech}.

\chapterpage{1 \enspace THE FIRST PART}{A chapter with a figure}

\Exact\ A tagged sentence. Figure~\ref{fig:stub} shows two dots and a link, and Chapter~2 has a
table. A mention of Section~II of a technical paper stays plain text on a post that has no technical
page to link to.

\begin{figure}[H]
\centering
\begin{tikzpicture}[x=1cm,y=1cm,font=\small]
\draw[flat,thick] (0,0) -- (2,0);
\fill[curl] (0,0) circle (0.08);
\fill[curl] (2,0) circle (0.08);
\node[text=ink] at (1,-0.4) {a link};
\end{tikzpicture}
\caption{Two dots and a link. \Sketch\ Drawn only so that the build has a figure to render.}
\label{fig:stub}
\end{figure}

\sub{A sub-heading}
A paragraph under it, with a number, $N=4$, and \emph{emphasis}.

\takeaway{The first box. The page keeps it as an aside.}

\chapterpage{2 \enspace THE SECOND PART}{A chapter with a table and a list}

\whyhere{A fixture needs one of each block.}

\sub{A sub-heading}
\Meas\ A table, centered:

\begin{center}
\renewcommand{\arraystretch}{1.2}
\begin{tabular}{lc}
\toprule
Thing & Count\\
\midrule
Dots & 2\\
Links & 1\\
\bottomrule
\end{tabular}
\end{center}

And a list:
\begin{itemize}\setlength{\itemsep}{1pt}
\item \Pub\ One item.
\item \Ours\ Another, which points back to Chapter~1.
\end{itemize}

\takeaway{The second box. There is no Chapter~7 here, so that mention is not a link.}

\par\Needspace{14\baselineskip}\vspace{1.4em}
{\color{ink}\LARGE\bfseries Continuing into the technical paper}\par
\vspace{0.6em}

A closing paragraph, which cites the technical paper once more~\cite{tech}.

\sub{A few translations}
\renewcommand{\arraystretch}{1.3}
\begin{tabularx}{\linewidth}{@{}>{\bfseries}p{0.24\linewidth}X@{}}
\toprule
In the paper & What to keep in mind\\
\midrule
$N$ & Number of dots.\\
Link & A line between two dots.\\
\bottomrule
\end{tabularx}

\renewcommand{\refname}{Sources}
\begin{thebibliography}{99}\small
% the first entry ends a line with a control space, which the build has to read as one
\bibitem{first} A.~Author, \emph{A first source}, J.\ Fixt.\
\textbf{1}, 1 (2000).
\bibitem{second} B.~Author, \emph{A second source}, \url{https://example.org/second}.
\bibitem{tech} C.~Author, \emph{The technical paper this fixture stands beside}, technical manuscript.
\end{thebibliography}

\end{document}
